% Ad Familiares 13.50 --- headnote
% ad-familiares-13-50, early January 44 BC, Rome

\textit{Cicero to Acilius, proconsul of Achaia, written from Rome
at the beginning of January 44 BC (the manuscript dateline:
\textit{Scr. Romae in. m. Ian. a. 710 (44)}). The beneficiary is
the same M'. Curius of Patrae who appears in 13.17 as commended
to Servius Sulpicius Rufus: a Roman businessman based at the chief
trading port of Achaia, an intimate of Atticus, and so the
particular favourite of Atticus's circle. Cicero leans on his
old Brundisium acquaintance with Acilius --- forged during the
unhappy months of his return from Pharsalus in 47 --- to extract
the broadest possible undertaking on Curius's behalf, that he be
kept (in the proverb of the leasehold contract) \textit{sartum et
tectum}, ``patched and roofed,'' against every kind of trouble.}
